% ═══════════════════════════════════════════════════════════════
%  INTRODUCTION GÉNÉRALE
% ═══════════════════════════════════════════════════════════════

\chapter*{Introduction Générale}
\addcontentsline{toc}{chapter}{Introduction Générale}

Soutenu par une urbanisation massive, le secteur des infrastructures et de l'immobilier camerounais connaît une croissance notable \hyperref[ref:africanvestor]{[7]}, pesant de manière significative dans le PIB du secteur tertiaire et des services \hyperref[ref:beac]{[17]}, avec une pénétration internet et mobile en progression soutenue \hyperref[ref:datareportal]{[8]}. Toutefois, cette dynamique est freinée par \textbf{une forte fragmentation de l'information : les annonces immobilières sont dispersées sur de nombreuses plateformes en ligne sans aucun mécanisme centralisé de vérification} \hyperref[ref:ndah]{[1]}. L'Institut National de la Statistique du Cameroun relève par ailleurs qu'une très large majorité des opérateurs intermédiaires évoluent dans le secteur informel \hyperref[ref:ins]{[18]}, ce qui entraîne des incohérences tarifaires massives limitant considérablement la transparence et la confiance des utilisateurs dans les outils existants \hyperref[ref:strataxis]{[2]}.

C'est dans ce contexte que se pose la question centrale de notre travail : \textbf{Comment réduire la fragmentation de l'information immobilière et offrir aux utilisateurs des annonces fiables via une application mobile et web ?} Cette question se décline en trois axes opérationnels portant sur les techniques de collecte automatique de données, la conception d'un algorithme de déduplication efficace et la mise à disposition des données via une API REST couplée à des interfaces mobiles et web.

Pour y répondre, nous avons poursuivi des objectifs précis et quantifiés dans le cadre d'un stage de six mois au sein de la startup Smart Service Hub (SSH), du 5 janvier au 30 juin 2026 : développer un pipeline de collecte couvrant au moins trois sources, concevoir un algorithme de déduplication capable d'identifier les annonces en double entre sources, garantir un temps de réponse de l'API inférieur à 500~ms, et livrer une application mobile et web fonctionnelle intégrant recherche, filtres et affichage cartographique.

Ce projet nous a motivé à la fois sur le plan personnel, car il répond à un besoin concret des utilisateurs camerounais, et sur le plan académique, car il mobilise des compétences clés de notre formation en développement logiciel, traitement de données et conception d'architectures logicielles multi-couches. Ce rapport s'organise en trois chapitres : le premier présente la structure d'accueil et son environnement technique ; le deuxième détaille la conception et le développement de notre plateforme ; le troisième évalue les résultats obtenus et propose un bilan des compétences acquises. Une conclusion générale clôturera ce travail.
